\documentclass[11pt,a4paper]{article}
\usepackage{fontspec}
\setmainfont{Latin Modern Roman}
\setmonofont{Latin Modern Mono}[Scale=MatchLowercase]
\usepackage{unicode-math}
\setmathfont{Latin Modern Math}
\usepackage{newunicodechar}
\newunicodechar{λ}{\ensuremath{\lambda}}
\newunicodechar{σ}{\ensuremath{\sigma}}
\newunicodechar{θ}{\ensuremath{\theta}}
\newunicodechar{≈}{\ensuremath{\approx}}
\newunicodechar{≥}{\ensuremath{\geq}}
\newunicodechar{≠}{\ensuremath{\neq}}
\newunicodechar{×}{\ensuremath{\times}}
\newunicodechar{↔}{\ensuremath{\leftrightarrow}}
\newunicodechar{→}{\ensuremath{\rightarrow}}
\newunicodechar{−}{\ensuremath{-}}
\newunicodechar{⁴}{\textsuperscript{4}}
\newunicodechar{⁵}{\textsuperscript{5}}
\newunicodechar{⁶}{\textsuperscript{6}}
\newunicodechar{⁷}{\textsuperscript{7}}
\newunicodechar{⁸}{\textsuperscript{8}}
\usepackage[a4paper,margin=1in]{geometry}
\usepackage{longtable,booktabs,array,calc}
\usepackage{pdflscape}
\usepackage{microtype}
\usepackage[hidelinks]{hyperref}
\setcounter{secnumdepth}{1}
\setcounter{tocdepth}{1}
\setlength{\emergencystretch}{3em}
\sloppy
\setlength{\parskip}{0.5\baselineskip}
\setlength{\parindent}{0pt}
\providecommand{\tightlist}{\setlength{\itemsep}{0pt}\setlength{\parskip}{0pt}}
\hyphenation{HNSCC MesenContext MesenCoder EMToolbox RESULT COMPLETION}
\begin{document}
{\Large\bfseries Supplementary Methods\par}\vspace{8pt}
{\bfseries EMToolbox uncovers biased co-option of embryonic EMT programs across cancers\par}\vspace{2pt}
Caleb C. Reagor, CO-AUTHORS, and Igor Adameyko\par\vspace{14pt}
{\bfseries MesenContext: Agent instruction set\par}\vspace{2pt}
MesenContext uses natural language to encode a decision tree and successive rounds of
human-agent knowledge transfer, indirectly steering AI agents during their autonomous analyses
of Curated Cancer Cell Atlas (3CA) datasets (Tyler \textit{et al.}, 2025). Each section of these methods was a separate specification file in \texttt{mesencontext/spec/}, reproduced below with cosmetic updates. Section 1 defines the analysis task, Sections 2--4 audit agent performance on steered and
autonomous analyses, and Section 5 logs all runs across 3CA datasets. The documented context was accumulated as the project progressed --- each agent
received every section completed before its run, including newly added subsections and ledger
updates. These
specification files were AI-generated with human steering to ensure scientific and methodological
accuracy.\par\vspace{12pt}
{\setlength{\parskip}{0pt}\tableofcontents}

\hypertarget{sec:1}{%
\section{Cell-state analysis: Task definition}\label{sec:1}}

\hypertarget{overview}{%
\subsection{0. Overview}\label{overview}}

Re-run a six-step single-cell analysis --- hand-built on Puram2017 (HNSCC) --- across every in-scope 3CA dataset, one per
subdirectory.

\textbf{Question:} Do malignant cells recapitulate embryonic-development
programs? Score three modules (Gast, pSom, Crest) on malignant cells,
fit a cell-state manifold / trajectory, and use three GO/NO-GO gates to
decide whether the expected pattern --- a Gast/pSom hub diverging into
two arms with cross-repression --- reproduces.

\textbf{Role:} This is \textbf{not} a fixed pipeline; nearly every
parameter is dataset-dependent. Act as the \textbf{analyst} --- make the
same \emph{kinds} of judgment calls made for Puram2017, grounded in the
data and biology. The Puram2017 notebooks are the \textbf{protocol}
(§1.3), not values to copy.

\textbf{Scope:} \texttt{CCCA\_summary.csv} (§1.4) is authoritative. The
on-disk folder set is a superset --- folders not in the summary are out
of scope; do not analyze them.

\hypertarget{biology}{%
\subsection{1. Biology}\label{biology}}

Hypothesis: Malignant cells recapitulate \textbf{embryonic-development
gene programs}; EMT mirrors developmental trajectories and/or
bifurcations. Modules (read the \texttt{hsapiens} column, drop
NaN/dupes):

\begin{itemize}
\tightlist
\item
  \textbf{Gast}: Gastrulation (\texttt{GSE162534.csv})
\item
  \textbf{pSom}: Posterior somitogenesis (\texttt{GSE229103.csv})
\item
  \textbf{Crest}: Neural crest (\texttt{rRNAModifications.csv} ---
  filename historical; contents are neural crest)
\item
  \textbf{union}: Control pool for \texttt{score\_genes}
  (\texttt{union.csv})
\end{itemize}

\textbf{Expected pattern (from Puram2017):} A \textbf{hub} where
Gast and pSom co-express \textbf{diverges} into a pSom-leaning arm and a
Gast-leaning arm. \textbf{Crest} is high at the arm tips (terminal
readout), enriched on the pSom arm. \textbf{Gast↔pSom is the
cross-repression axis} --- opposite directions out of the hub; mutual
exclusivity is the bifurcation signature. Crest is a terminal readout,
\textbf{not} on the cross-repression axis. Hub module correlation in
Puram ≈ \textbf{−0.49}. Module→arm assignment is \textbf{data-driven}
--- read it off the scores, not off subtype names.

\textbf{Not every dataset bifurcates.} In the \textbf{non-branching
case}, Gast and pSom \textbf{partially or fully overlap}, with Crest as
an inverse/terminal gradient --- the programs are present but not
resolved into two competing arms, though Gast and pSom may still display
some cross-repression. This is a real outcome, not a failure: it
\textbf{passes Gate 2} (\emph{i.e.}, GO), has no bifurcation for Gate 3
(\emph{i.e.}, NA), and is reported as such --- distinct from a NO-GO
where the pattern is absent.

\hypertarget{the-three-gono-go-gates-the-crux}{%
\subsection{2. The three GO/NO-GO gates (the crux)}\label{the-three-gono-go-gates-the-crux}}

Gates are \textbf{scientific} (is the phenomenon present?), not
parameter-quality checks. Parameter tuning is the agent's means to
answer them honestly --- never something the researcher gates on; judged
visually.

\begin{itemize}
\tightlist
\item
  \textbf{Gate 1 (01→02):} After Harmony on \texttt{sample} and DR, is the
  malignant structure both \textbf{mixed} (sample-integrated) and
  \textbf{continuous} (connected gradient)? NO-GO if it clusters by
  sample or breaks into disconnected islands --- don't force a
  trajectory (importance level = HIGH).
\item
  \textbf{Gate 2 (03→04):} Does the embryonic module pattern appear ---
  a Gast/pSom hub diverging into two arms, Crest high at tips (ideally
  pSom arm)? The \textbf{non-branching case} (Gast/pSom overlap, Crest
  an inverse gradient; §1.1) still \textbf{passes} (\emph{i.e.}, GO) but
  has no bifurcation for Gate 3. \textbf{Flag} (pass anyway) if Crest
  enriches the \textbf{Gast} arm --- an inversion.
\item
  \textbf{Gate 3 (04→05, BRANCHING ONLY):} Do tree nodes show
  \textbf{mutual Gast↔pSom negative correlation} (competitive
  cross-repression)? Crest not tested. No anti-correlation → NO-GO on
  step 05.
\end{itemize}

\textbf{The EMT-TF expression step is unconditional} --- always produced
as the final notebook, outside the gate chain (numbered 06 only when all
gates pass; renumbered after an earlier NO-GO --- see §1.6).

\textbf{Linchpin: identify the HUB} (\emph{e.g.}, Puram's ``Mixed'') ---
the Gast/pSom-overlap convergence population that 04--05 run on and Gate
3 tests. Identify by \textbf{role} (module-convergence zone at the
trajectory branch point/root), not by name, and confirm it's the
topological convergence point.

\hypertarget{the-6-step-protocol-source-of-truth-puram2017-notebooks}{%
\subsection{3. The 6-step protocol (source of truth: Puram2017 notebooks)}\label{the-6-step-protocol-source-of-truth-puram2017-notebooks}}

\texttt{.../CCCA/Puram2017\_Head-and-Neck/\{01..06\}\_*.ipynb}
(READ-ONLY). Parameter values in reference notebooks are
\textbf{illustrative, not defaults to copy} --- and the
\textbf{in-notebook comments are illustrative of the analytical judgment
calls} the researcher made: the reasoning to emulate rather than
transcribe. I/O per step in §1.6.

\begin{enumerate}
\def\labelenumi{\arabic{enumi}.}
\tightlist
\item
  \textbf{01\_preprocessing}: Load malignant cells → drop
  under-populated samples (per-dataset min-cell threshold) → filter
  genes → overdispersed selection → PCA → Harmony on \texttt{sample} →
  neighbors → UMAP → Leiden → drop detached clusters → re-cluster →
  annotate subtypes. \textbf{Annotation always uses both data-derived
  marker expression and the original publication's reported
  subtypes/markers} (cited in the log) --- never one alone.
\item
  \textbf{02\_manifold}: ScFates PPT tree (branching) / curve
  (non-branching) → biologically-motivated root → pseudotime → rename
  milestones.
\item
  \textbf{03\_signatures}: \texttt{score\_genes} for Gast/pSom/Crest
  with the union control pool → signature UMAPs.
\item
  \textbf{04\_correlation}: Gast XOR pSom genes, hub cells →
  pseudotime-windowed module Pearson → correlation UMAP.
\item
  \textbf{05\_modules}: Gast OR pSom genes, hub cells → per-gene
  module correlation → mutual-correlation scatter (reports \emph{r}).
\item
  \textbf{06\_expression}: EMT-TF expression overlay; most
  informative TFs for that dataset.
\end{enumerate}

Per-dataset notebooks are numbered \textbf{sequentially} --- the
expression step is always the \emph{last} notebook, so its number
depends on how far the analysis ran (§1.6).

\hypertarget{data-code-assets-source-of-truth}{%
\subsection{4. Data \& code assets (source of truth)}\label{data-code-assets-source-of-truth}}

\begin{itemize}
\tightlist
\item
  \textbf{Loader} \texttt{.../mesencontext/scripts/load.py} ---
  \texttt{load\_CCCA\_adata(datadir)} reads the matrix/gene/cell files
  in \texttt{datadir}, auto-detects normalization from the matrix
  filename (TPM/normalized pass through; counts normalized), always
  \texttt{log1p}, sets \texttt{celltype} (default `Malignant').
  \textbf{Point it at a bottom-level (leaf) dataset directory}, as
  \texttt{merge.py} does --- it does not walk nested subdirs.
  \texttt{merge.py} also defines the folder-naming convention mapping
  summary rows → on-disk folders.
\item
  \textbf{Scope list} \texttt{.../mesencontext/claude/CCCA\_summary.csv}
  --- TAB-sep;
  \texttt{Title,\ Category,\ Disease,\ Technology,\ \ \ Samples,\ Cells};
  one row per in-scope dataset. \textbf{Authoritative.}
\item
  \textbf{Input data} \texttt{.../data/unzip/CCCA/} --- one folder per
  dataset (superset of scope). Layout varies (some datasets nest by
  technology; expression is a mix of counts and TPM). Resolve each
  dataset to its leaf matrix directory before loading; don't hardcode
  format counts.
\item
  \textbf{Gene sets} \texttt{.../data/features/biomart/} --- the four
  CSVs above.
\item
  \textbf{Environment}: The existing \texttt{scFates} conda env (§1.5).
  Runable by setting \texttt{R\_HOME} path in bash (per notebooks).
\end{itemize}

\hypertarget{directory-access-environment}{%
\subsection{5. Directory access \& environment}\label{directory-access-environment}}

Grant \textbf{narrowly, one at a time}; grants apply to \textbf{new
kernels only} --- restart after each.

{\small
\begin{longtable}[]{@{}
  >{\raggedright\arraybackslash}p{(\columnwidth - 4\tabcolsep) * \real{0.540}}
  >{\raggedright\arraybackslash}p{(\columnwidth - 4\tabcolsep) * \real{0.110}}
  >{\raggedright\arraybackslash}p{(\columnwidth - 4\tabcolsep) * \real{0.350}}@{}}
\toprule\noalign{}
\begin{minipage}[b]{\linewidth}\raggedright
Directory
\end{minipage} & \begin{minipage}[b]{\linewidth}\raggedright
Access
\end{minipage} & \begin{minipage}[b]{\linewidth}\raggedright
Role
\end{minipage} \\
\midrule\noalign{}
\endhead
\bottomrule\noalign{}
\endlastfoot
\texttt{.../data/unzip/CCCA/} & \textbf{ro} & input dataset folders \\
\texttt{.../mesencontext/scripts/} & \textbf{ro} & \texttt{load.py},
\texttt{merge.py} \\
\texttt{.../mesencontext/Puram2017\_Head-and-Neck/} & \textbf{ro} &
ground-truth reference notebooks 01--06 \\
\texttt{.../data/features/biomart/} & \textbf{ro} & gene sets \\
\texttt{.../mesencontext/claude/} & \textbf{rw} & notebooks, summary-CSV
(copy), cumulative log \\
\texttt{.../data/cancer/vertebrate/CCCA/claude/} & \textbf{rw} & output
h5ads \\
\texttt{.../figures/mesencontext/claude/} & \textbf{rw} & output figures (with \texttt{colorbars/}) \\
\end{longtable}
}

Not granted: \texttt{EMToolbox} root, \texttt{data} root --- hence the
summary CSV lives in the rw \texttt{claude/}.

\textbf{Environment:} Use the existing \textbf{\texttt{scFates}} conda
env --- \textbf{no recreation needed}. (Stack: python 3.11, r-base 4.3 with r-mgcv, rpy2, scanpy/scFates/harmonypy/leidenalg/anndata; scFates PPT →
rpy2 → r-mgcv, and notebook 01 sets \texttt{R\_HOME} to the conda
prefix.)

\hypertarget{deliverables-outputs}{%
\subsection{6. Deliverables \& outputs}\label{deliverables-outputs}}

\textbf{Layout:} One dataset per subdir, partitioned by
\texttt{Category}, replicated across three roots (notebooks, h5ads, figures and \texttt{colorbars/}):
\texttt{.../claude/\textless{}Category\textgreater{}/\textless{}DatasetName\textgreater{}/NN\_*.ipynb}
(\texttt{\textless{}DatasetName\textgreater{}} = folder name minus
\texttt{Data\_}). \textbf{Relative paths:} The
\texttt{\textless{}Category\textgreater{}/} partition puts notebooks two
levels deeper than Puram --- adapt the path logic; notebooks must run in
place. \textbf{Keep rw-directory outputs tidy and as brief as possible}
--- they live on the researcher's disk; don't accumulate scratch or
intermediate files. These deliverables and instructions are
\textbf{BINDING}.

{\small
\begin{longtable}[]{@{}
  >{\raggedright\arraybackslash}p{(\columnwidth - 6\tabcolsep) * \real{0.140}}
  >{\raggedright\arraybackslash}p{(\columnwidth - 6\tabcolsep) * \real{0.230}}
  >{\raggedright\arraybackslash}p{(\columnwidth - 6\tabcolsep) * \real{0.230}}
  >{\raggedright\arraybackslash}p{(\columnwidth - 6\tabcolsep) * \real{0.400}}@{}}
\toprule\noalign{}
\begin{minipage}[b]{\linewidth}\raggedright
Step
\end{minipage} & \begin{minipage}[b]{\linewidth}\raggedright
Reads
\end{minipage} & \begin{minipage}[b]{\linewidth}\raggedright
h5ad
\end{minipage} & \begin{minipage}[b]{\linewidth}\raggedright
Figures
\end{minipage} \\
\midrule\noalign{}
\endhead
\bottomrule\noalign{}
\endlastfoot
preprocessing & raw load & \texttt{processed.h5ad} &
\texttt{samples.pdf} \\
manifold & \texttt{processed.h5ad} & \texttt{manifold.h5ad} &
\texttt{manifold.pdf} \\
signatures & \texttt{manifold.h5ad} & \texttt{signatures.h5ad} &
\texttt{\{Gast,pSom,Crest\}\_signature.pdf}, colorbars \\
correlation & \texttt{signatures.h5ad} & --- & \texttt{correlation.pdf},
\texttt{colorbars/R\_colorbar.pdf} \\
modules & \texttt{signatures.h5ad} & --- & \texttt{modules.pdf} \\
expression & \texttt{processed.h5ad} & --- & \texttt{expression.pdf} \\
\end{longtable}
}

\textbf{Stop at a NO-GO: do not create downstream notebooks} (not
stubbed, not emptied). Notebook numbers are \textbf{sequential}, so the
unconditional expression step is always the \emph{last} notebook,
renumbered to follow whatever preceded it:

\begin{itemize}
\tightlist
\item
  Gate 1 NO-GO → preprocessing (01) and expression (\textbf{02})
\item
  Gate 2 NO-GO → 01--03 and expression (\textbf{04})
\item
  Gate 3 NO-GO → 01--04 and expression (\textbf{05})
\item
  all gates pass → full 01--06 (expression = 06)
\end{itemize}

\textbf{Notebooks should contain alternating code and figures} ---
matching the working style of the ground-truth Puram2017 notebooks.
These deliverables should reflect the \emph{exact} analysis and results
that inform the GO/NO-GO decisions. Any inconsistencies between the
notebooks and scientific analyses will \textbf{fully invalidate} the
results --- so treat these notebooks as \emph{critical}
assets/deliverables.

\textbf{Every dataset also produces, regardless of how far it ran:}

\begin{itemize}
\tightlist
\item
  \textbf{\texttt{summary.pdf}}: A \textbf{punchy} one-page executive
  summary (same register as this doc): dataset in one line (cancer type,
  platform, malignant-cell count); key parameter choices and rationale
  (min-cell threshold, root node, geometry); \textbf{the gate outcomes
  with reasoning for each GO/NO-GO}; hub identity and r, where
  applicable; a one-line verdict on whether the pattern reproduces.
  \textbf{Embeds the step figures} (the produced \texttt{.pdf}s) so the
  decisions are visible at a glance. Format document to match the
  Puram2017 calibration-step output --- this is \textbf{BINDING}.
\item
  \textbf{one row in \texttt{.../mesencontext/claude/analysis\_log.csv}}
  (TAB-sep) --- the six \texttt{CCCA\_summary.csv} columns verbatim
  (joins 1:1 to the master), then provenance (join key, cell counts
  through filtering, expression type, key embedding/trajectory params,
  clusters removed, subtypes and annotation source, geometry) and the gate
  readout (\texttt{gate1/2/3}, Crest-arm flag, hub r, EMT-TFs shown,
  status, notes).
\end{itemize}

\hypertarget{sec:2}{%
\section{Calibration audit: Puram2017 HNSCC}\label{sec:2}}

The autonomous analyst re-ran the six-step protocol on \textbf{Puram2017}, the dataset it was hand-built on. Reached step 06, \textbf{all
three gates GO}. Same data, known answer --- so the real question is not
\emph{did it pass} but \emph{which passes reflect a decision that will
survive contact with data the agent has never seen}.

\textbf{Verdict.} The pattern reproduces: Mixed hub confirmed,
cross-repression \textbf{r = −0.47} (anchor −0.49), hub-and-spoke
geometry intact. Cleared to scale, with the graded watch-list below.

\hypertarget{automation-criteria-the-bar-a-decision-must-clear-to-run-unattended}{%
\subsection{Automation criteria: The bar a decision must clear to run unattended}\label{automation-criteria-the-bar-a-decision-must-clear-to-run-unattended}}

Every dataset-dependent choice must be \emph{automated} for the 60+
dataset run. A choice is safely automated only when it clears four
tests; anything short is a liability that a passing gate on the training
data (\emph{i.e.}, Puram2017) will hide:

\begin{itemize}
\tightlist
\item
  \\textbf{C1: Derived, not inherited.} Computed from the dataset in
  hand, not copied from the anchor run.
\item
  \\textbf{C2: Shown, not asserted.} The derivation is demonstrated in
  the notebook, not just claimed in a comment (importance =
  \textbf{HIGH}).
\item
  \\textbf{C3: Assumptions stated and met.} The rule's structural
  premises (one hub, one leftover cluster, a clean count gap, a clean
  bifurcation) are made explicit and hold on the target data.
  Circularity avoided.
\item
  \\textbf{C4: Portable across scope.} Generalizes over the atlas's
  actual range: cancer types, platforms (10x vs SmartSeq2), and cell
  counts spanning \textasciitilde2k to \textasciitilde90k.
\end{itemize}

The fixed \texttt{Puram2017} skeleton is out of scope here: where the
protocol prescribes a value, the agent simply reused it (malignant load
2,539/17 pt, threshold ≥100 → 2,472, overdispersion 4,804,
Harmony/neighbors/UMAP, PPT params, union-pool scoring --- down to gene
recovery 78/92/103). The criteria apply only to the \emph{tunable}
choices below.

\hypertarget{five-decisions-re-derived}{%
\subsection{Five decisions re-derived}\label{five-decisions-re-derived}}

All five are genuinely \emph{derived}, not inherited (C1) --- the real
achievement. The failures cluster instead on unstated assumptions (C3)
and scope (C4), exactly what a training-data calibration cannot expose.
Each below: Method, then assessment.

\begin{enumerate}
\def\labelenumi{\arabic{enumi}.}
\tightlist
\item
  \textbf{Sample threshold} (≥100 → 2,472). \emph{Method:} Sort
  per-sample malignant cell counts, read the 23→100 jump as an
  inflection, keep samples above it. \emph{Assessment:} Depth-invariant
  in principle, but no gap-finding is actually run (C2) --- the
  ``inflection'' is just the anchor's ≥100 written as a constant; smooth
  count distributions have no gap, leaving the rule ill-defined (C3).
\item
  \textbf{Island removal} (−64 → 2,408, vs reference's hand-listed
  Leiden −116). \emph{Method:} Build a UMAP radius graph at 6× the
  median nearest-neighbour distance, take connected components, keep the
  largest. \emph{Assessment:} Automated and strictly conservative ---
  strips only disconnected specks --- but the 6× radius is an undeclared
  constant, and ``keep the largest component'' assumes noise detaches
  cleanly (C3): on a dense 60k-cell 10x manifold debris stays attached,
  or a real second lineage reads as a component and is deleted (C4).
  Biggest numerical gap (\textbf{+52 cells}).
\item
  \textbf{Subtype annotation} (Basal/Classical/Atypical/Mixed).
  \emph{Method:} Leiden at resolution 0.15, then assign each label to
  the cluster with the highest mean marker-program score (Basal
  COL17A1/KRT14/TP63; Classical NRF2 program; Atypical XAGE1 family),
  greedily, remainder = Mixed. \emph{Assessment:} Robust to Leiden
  renumbering --- no per-dataset re-mapping --- but the panels are
  HNSCC-specific and won't transfer across cancer types (C4),
  ``remainder = Mixed'' assumes exactly one leftover cluster (C3), and
  resolution 0.15 is hardcoded.
\item
  \textbf{Root node} (node 85 vs manual 185; same biology).
  \emph{Method:} Fit the PPT tree, sum the soft-assignment
  responsibility of Mixed-hub cells over principal nodes, root at the
  argmax. \emph{Assessment:} Removes the last human step --- the change
  that makes unattended scale-up viable --- but it inherits every
  upstream error and assumes exactly one convergence point (C3). Most
  fragile step for unseen datasets.
\item
  \textbf{EMT-TF readout} (ZEB1 and PRRX1 vs hand-picked SNAI1 and ZEB1).
  \emph{Method:} From an 8-TF candidate panel, score each by
  arm-specificity (top-arm mean / total arm expression) among expressed
  factors, take the best TF for two \emph{different} arms.
  \emph{Assessment:} The procedure generalizes and correctly rejects the
  high-but-nonspecific SNAI2, but it shows different genes per dataset
  --- so the panel is no longer comparable across the atlas --- and
  presumes a clean bifurcation (C3). Only panel not anchor-matched.
\end{enumerate}

\hypertarget{overfitting-three-choices-that-fail-c1-outright}{%
\subsection{Overfitting: Three choices that fail C1 outright}\label{overfitting-three-choices-that-fail-c1-outright}}

Three continuous knobs were copied verbatim, rationale asserted but
never derived (fail C1 and C2): \textbf{PPT (Nodes=200, λ=2000, σ=0.5)}
--- all scale with cell count, so 200 nodes over 2,408 cells ≠ over
60,000; \textbf{n\_pcs=40 / spread=1.3} --- unexamined for
larger/sparser sets; \textbf{Gate-3 bar r \textless{} −0.2} --- never
tested against a null, and it never bit because −0.47 sits well below
it. On the training data an inherited constant passes the gate
\emph{because} it was fitted here --- the gate certifies nothing about
generalization. This is what C1 exists to catch.

\hypertarget{gates-ledger}{%
\subsection{Gates \& ledger}\label{gates-ledger}}

G1 (mixing/continuity) GO. G2 (module geography) GO --- Crest peaks at
Basal tip (0.43 vs 0.27 ≈ 0.24), Gast highest in Classical, pSom flat.
G3 (hub cross-repression) GO --- \textbf{r = −0.47}, within 0.02 of the
anchor.

{\small
\begin{longtable}[]{@{}
  >{\raggedright\arraybackslash}p{(\columnwidth - 6\tabcolsep) * \real{0.300}}
  >{\raggedright\arraybackslash}p{(\columnwidth - 6\tabcolsep) * \real{0.240}}
  >{\raggedright\arraybackslash}p{(\columnwidth - 6\tabcolsep) * \real{0.220}}
  >{\raggedright\arraybackslash}p{(\columnwidth - 6\tabcolsep) * \real{0.240}}@{}}
\toprule\noalign{}
\begin{minipage}[b]{\linewidth}\raggedright
Step
\end{minipage} & \begin{minipage}[b]{\linewidth}\raggedright
Reference
\end{minipage} & \begin{minipage}[b]{\linewidth}\raggedright
Agent
\end{minipage} & \begin{minipage}[b]{\linewidth}\raggedright
Divergence
\end{minipage} \\
\midrule\noalign{}
\endhead
\bottomrule\noalign{}
\endlastfoot
Load / threshold / overdispersion & 2,539·2,472·4,804 &
2,539·2,472·4,804 & identical \\
Neighbors / UMAP & n\_pcs=40, spread=1.3 & same & identical (inherited,
fails C1) \\
\textbf{Island removal} & Leiden 1.5 → 2,356 & radius-graph → 2,408 &
\textbf{method · +52} \\
\textbf{Subtype annotation} & manual ID→label & marker-rule auto &
\textbf{method (same 4 labels)} \\
Trajectory (PPT) & Nodes=200 λ=2000 σ=0.5 & same & identical (inherited,
fails C1) \\
\textbf{Root node} & 185 (manual) & 85 (argmax Mixed) & \textbf{method
(both = hub)} \\
\textbf{Gate 3 hub r} & −0.49 (n=476) & \textbf{−0.47 (n=466)} &
\textbf{match} \\
\textbf{EMT-TF pair} & SNAI1/ZEB1 (manual) & ZEB1/PRRX1 (auto) &
\textbf{readout diverges} \\
\end{longtable}
}

\textbf{Bottom line.} Real automation, not mimicry: all five tunable
decisions pass C1. But none was stress-tested on C3/C4, and three
hyperparameters fail C1 entirely. Cleared to scale --- on the first
dataset where the answer is unknown, re-derive the three inherited
constants and watch the C3/C4 assumptions (one hub, one leftover
cluster, a clean bifurcation, a clean count gap) for the first case that
violates them.

\hypertarget{sec:3}{%
\section{Evaluation audit: Kürten2021 HNSCC}\label{sec:3}}

The analyst ran the six-step protocol on \textbf{Kürten2021}, the first
dataset whose answer was not known in advance. Reached the expression
step, \textbf{G1 GO / G2 GO / G3 NA} --- the non-branching case (Section
1 §1.1). An expert steered step-by-step and a fresh-context reviewer
checked claims. The question is not \emph{did it pass} but \emph{which
passes came from the agent's own reasoning} --- the gap a
fully-autonomous run must close.

\textbf{Verdict.} The pattern reproduces: single embryonic axis, root
Gast/pSom-high → terminus Crest-high, EMT-TFs \textbf{FOXC1 → SNAI2} ---
a strict linear subset of Puram. But at nearly every gate the agent's
first call was wrong, corrected by the expert or reviewer. Cleared as a
\emph{supervised} run; not as evidence of autonomy.

\hypertarget{six-forks}{%
\subsection{Six forks}\label{six-forks}}

\begin{enumerate}
\def\labelenumi{\arabic{enumi}.}
\item
  \textbf{Gate 1.} The agent passed G1 on a mixing scalar and a
  connected-components test; the expert had to redirect attention to the
  embedding itself, which showed a per-patient archipelago the metrics
  had scored as integrated. \emph{Rule:} Judge G1 on the regenerated
  embedding across an integration-strength sweep, with a
  per-neighborhood metric (iLISI); never on a scalar alone.
\item
  \textbf{Lobe identity.} The integrated manifold resolved into two
  lobes and the agent moved to fit one curve across both; the expert
  required each lobe's biological identity first. Only one lobe
  reproduced the embryonic signature relationships (Gast/pSom co-vary,
  both oppose Crest); the other was a cycling compartment with Gast
  confounded by proliferation. \emph{Rule:} A manifold resolves into
  sub-regions --- settle their biology (signatures and cell cycle)
  before the G1 GO, and scope the trajectory to the region that
  reproduces the reference relationships.
\item
  \textbf{Root \& topology.} The agent rooted at the Crest tip, treating
  neural crest as the progenitor state; developmentally Crest is
  \emph{late}, so the root belongs at the Gast/pSom tip. A non-branching
  topology is fit with scFates \texttt{curve}, not \texttt{tree}.
  \emph{Rule:} Root from the signatures' developmental order; match the
  scFates method to the topology.
\item
  \textbf{Continuity.} The agent certified the curve as continuous using
  a metric built from UMAP neighbors, blind by construction to a fold in
  that same UMAP; the reviewer caught the fold. It arose from
  re-embedding a linear compartment --- an elongated strip a fresh UMAP
  folds into a U. \emph{Rule:} Do not re-embed a linear compartment, and
  do not measure continuity with a metric that reads its neighbors off
  the embedding being tested.
\item
  \textbf{EMT-TF readout.} The canonical Snail/Zeb/Twist panel yielded
  only a terminus factor; FOXC1, the root-end factor, surfaced only once
  the panel was widened. \emph{Rule:} Score EMT-TFs on this dataset
  gated on detectable expression, and include other mesenchymal families
  such as FOXC and PRRX, not only Snail/Zeb/Twist.
\item
  \textbf{Deliverables.} With G3 = NA, the agent began the
  branching-only \texttt{04\_correlation} / \texttt{05\_modules}
  notebooks, which assert structure not found. \emph{Rule:} Gate
  outcomes decide which notebooks exist; an NA or NO-GO means the
  corresponding deliverables are not produced (§1.6).
\end{enumerate}

\hypertarget{three-lessons-for-an-autonomous-run}{%
\subsection{Three lessons for an autonomous run}\label{three-lessons-for-an-autonomous-run}}

\begin{enumerate}
\def\labelenumi{\arabic{enumi}.}
\item
  \textbf{Biology before geometry.} The scientifically decisive forks
  were the same failure --- a geometric result read as biology without
  the biological check. That check must run and appear in the notebook
  before the GO: G1 against the regenerated embedding, each sub-region's
  identity before the curve, continuity off an independent embedding.
  The gate stays binary --- the agent proposed ``Conditional GO'' twice,
  and it is not an outcome.
\item
  \textbf{Check against the spec, not memory.} The run's deliverable
  errors were caught by re-reading the binding spec against the
  notebooks rather than trusting recall --- the renumbering rule, the
  deliverable set, and the summary format were each re-verified against
  the spec text before finalizing. An autonomous run should re-read the
  relevant spec section at each gate and before shipping, and check its
  outputs against the binding text (importance = \textbf{HIGH}). An
  agent must ALWAYS re-read the full spec and context before passing gates
  --- this is now \textbf{BINDING}.
\item
  \textbf{Deliverables must be reproducible and consistent.} Each
  notebook must run end-to-end from raw and reach the reviewed result,
  and the values it produces must match the summary and log that report
  them. Where a step is stochastic (Harmony's KMeans is not
  seed-deterministic), state the tolerance so a benign shift is not read
  as a changed decision.
\end{enumerate}

\hypertarget{bottom-line}{%
\subsection{Bottom line}\label{bottom-line}}

The protocol transferred; the judgment did not. The derivation machinery
ran correctly on unseen data (threshold by gap-finding, \texttt{n\_pcs}
off the scree, detection-gated signatures), but the decisive calls ---
which lobe is the manifold, where the trajectory roots, which
transcription factors to read --- the agent got wrong on first pass, and
the expert corrected each by putting biology ahead of geometry. A fully
autonomous run makes those calls without the expert; the three lessons
above are what that takes. The verifier is the readiest to automate ---
the reviewer here was already a fresh-context agent --- and is the
safeguard that caught the two claims the agent made after asserting the
rule against them.

\hypertarget{action-items}{%
\subsection{Action items}\label{action-items}}

To safeguard the scientific accuracy of autonomous analyses and maintain
FULL compliance with the task definition and calibration/evaluation
audits, agents must \textbf{ALWAYS} re-read the \textbf{FULL} spec
before passing a gate, planning an analysis, or writing any deliverable
--- \textbf{THIS IS BINDING}.

\begin{itemize}
\tightlist
\item
  \textbf{FAILURE TO REGULARLY REVISIT THE FULL SPEC FILES WILL RESULT
  IN COMPOUNDING ERRORS AND INCORRECT TASK COMPLETION}
\end{itemize}

\hypertarget{sec:4}{%
\section{Testing audit: Autonomous runs}\label{sec:4}}

Each attempt on the \textbf{testing} datasets is a fully-autonomous run
--- the analysis is driven by the spec files alone, with no expert
steering. A fresh-context auditor then reviews the run, guided by the
expert. One entry per attempt below records the decisive points only;
read alongside the three spec files. Datasets are sectioned below,
attempts ordered within each.

\begin{center}\rule{0.5\linewidth}{0.5pt}\end{center}

\hypertarget{chen-2020-npc-10x}{%
\subsection{Chen2020 (NPC, 10x)}\label{chen-2020-npc-10x}}

Nasopharyngeal carcinoma, 10x, \textasciitilde3.7k malignant cells ---
head-and-neck, same tissue family as the Puram/Kürten ground-truth runs.

\hypertarget{attempt-1-gate-1-no-go}{%
\subsubsection{Attempt 1 · Gate 1 NO-GO}\label{attempt-1-gate-1-no-go}}

\begin{itemize}
\tightlist
\item
  \textbf{Fail: \texttt{n\_pcs=6} under-powers the manifold.}
  Chord-elbow over-cut the inherited 40; too few PCs force a
  patient-segregated read, so the NO-GO is likely an embedding artifact.
  (Elbow biases low here; the Puram-40 / Kürten-20 precedents are the
  scale to sanity-check against.)
\item
  \textbf{Fail: \texttt{summary.pdf} wrong layout.} Verdict-prose and stacked step-PDFs, not the reference gate-ledger → figure-grid
  one-pager.
\item
  \textbf{Held.} Spec re-read at each gate/deliverable; threshold, theta
  sweep, EMT-TF panel, renumbering all sound.
\item
  \textbf{Watch.} Gate call stayed metric-led --- embedding read as
  corroboration, not lead.
\end{itemize}

\hypertarget{attempt-2-gate-1-no-go}{%
\subsubsection{Attempt 2 · Gate 1 NO-GO}\label{attempt-2-gate-1-no-go}}

\begin{itemize}
\tightlist
\item
  \textbf{Fail: Integration strength too low.} Settled
  \texttt{theta=2} and swept only \texttt{theta\ 0→3}, dismissing
  \texttt{theta=3} as over-integration without a null test to justify
  the cutoff. The 10x precedent (Kürten) reached \texttt{theta=10}, so
  stronger integration was left untested before the NO-GO. The sweep
  covered \texttt{n\_pcs} well (attempt-1's issue) but not integration
  strength.
\item
  \textbf{Fail: Figure deliverables diverge from the reference.}
  \texttt{samples.pdf} is a bespoke 3-panel diagnostic (patient UMAP, cell-cycle and \texttt{n\_pcs×theta} sweep) built to argue the audit
  narrative, not the reference's single sample-colored manifold UMAP;
  and \texttt{summary.pdf} was written to the figures root, not the
  notebook dir where the reference keeps it.
\item
  \textbf{Held.} Both attempt-1 defects fixed: \texttt{n\_pcs=20}
  (matches the Kürten/Puram scale) and the \texttt{summary.pdf} layout
  now follows the reference. Gate 1 led with the embedding. Spec
  re-read/certified throughout.
\end{itemize}

\hypertarget{attempt-3-gate-1-go-gate-2-no-go-gate-3-na}{%
\subsubsection{Attempt 3 · Gate 1 GO / Gate 2 NO-GO / Gate 3
NA}\label{attempt-3-gate-1-go-gate-2-no-go-gate-3-na}}

\begin{itemize}
\tightlist
\item
  \textbf{Fail: Trajectory mis-specified.} Fit a branching
  \texttt{tree} on a manifold it had itself called non-branching (should
  be a \texttt{curve}), and rooted at a proliferation state rather than
  the developmental origin the signatures imply --- the same root error
  flagged on Kürten.
\item
  \textbf{Fail: Signature scores not smoothed.} Read geography and
  the Gast/pSom correlation from raw per-cell scores; the Kürten 10x
  reference graph-smooths the scores (knn-connectivity averaging in
  02\_manifold) before reading geography, so the raw read is
  noise-dominated here.
\item
  \textbf{Fail: EMT-TF panel degenerate.} The display cutoff was
  unset, so the color scale is outlier-dominated and washed out; and VIM
  was shown as an EMT-TF, but vimentin is a mesenchymal marker, not a
  transcription factor.
\item
  \textbf{Fail: Island removal incomplete.} Its metric was
  degenerate, so detached islands were left in the manifold and Gate 1
  was still called continuous.
\item
  \textbf{Held.} Deliverables now correct in format and location. Traced
  the two prior Gate-1 NO-GO calls to an \texttt{sc.pp.scale} step
  before PCA absent from the reference protocol --- a correct diagnosis,
  and why the chain finally ran past Gate 1. Spec re-read/certified
  throughout.
\end{itemize}

\hypertarget{attempt-4-stopped-pre-analysis-no-certification}{%
\subsubsection{Attempt 4 · Stopped pre-analysis (no
certification)}\label{attempt-4-stopped-pre-analysis-no-certification}}

\begin{itemize}
\tightlist
\item
  \textbf{Fail: No certification, immediate drift.} Never certified
  the binding re-read (it had slipped behind a context fold); drifted
  into building notebook-execution infrastructure and ran code before
  any gate. Stopped pre-analysis, then cleaned up correctly.
\end{itemize}

\hypertarget{attempt-5-gate-1-go-gate-2-no-go-gate-3-na}{%
\subsubsection{Attempt 5 · Gate 1 GO / Gate 2 NO-GO / Gate 3
NA}\label{attempt-5-gate-1-go-gate-2-no-go-gate-3-na}}

\begin{itemize}
\tightlist
\item
  \textbf{Fail: Linear curve fit to a visibly tree-shaped manifold.}
  Called the clearly Y-shaped manifold ``robustly non-branching'' on a
  PPT metric (0 branch nodes across 12 param combos) and fit one curve
  through it. The inverse of attempt 3's tree-on-a-strip; same cause ---
  topology read from the metric, not the embedding.
\item
  \textbf{Fail: Figure deliverables still diverge from the reference
  layouts, despite the spec re-read.} Re-reading the spec prose did not
  become a check of the outputs against the binding reference figures
  themselves.
\item
  \textbf{Fail: Gate 2 built toward GO, then reversed to NO-GO on a
  scalar comparison.} The module geography was legible in the signature
  embeddings; the call was driven by the scalar comparison instead.
\item
  \textbf{Held.} Binding re-read certified at start and at the gate
  (attempt-4 lesson landed); deliverables in the correct locations; no
  \texttt{sc.pp.scale}; \texttt{n\_pcs≈19}, theta swept.
\end{itemize}

\hypertarget{attempt-6-gate-1-go-gate-2-go-non-branching-gate-3-na}{%
\subsubsection{Attempt 6 · Gate 1 GO / Gate 2 GO (non-branching) / Gate
3 NA}\label{attempt-6-gate-1-go-gate-2-go-non-branching-gate-3-na}}

\begin{itemize}
\tightlist
\item
  \textbf{Fail: Folded curve, never caught by any post-hoc
  visualization.} The committed curve W-folds back through the cloud;
  \texttt{manifold.pdf} draws pseudotime color but not the curve edges,
  so topology was judged on signature↔pseudotime correlations, blind to
  the fold (the Kürten failure mode; Section 3). Surfaced only when
  the expert asked for a curve render in the debrief.
\item
  \textbf{Fail: Figure deliverables again diverge; format, content,
  and tidiness are BINDING.} manifold/samples retitled with milestone
  and legend labels overlaid on the cloud, not the reference's single panel and external colorbar; an extra \texttt{gate2\_diagnostics.pdf}
  left in the figures dir, outside the §1.6 set (scratch, as in
  attempt 2).
\item
  \textbf{Held.} Certified the binding re-read; reached Gate 2 GO
  (non-branching) matching Kürten; overturned its provisional tree→curve
  on biology (Gast/pSom co-vary, Crest inverse), not geometry.
\end{itemize}

\hypertarget{attempt-7-gate-1-go-gate-2-go-non-branching-gate-3-na}{%
\subsubsection{Attempt 7 · Gate 1 GO / Gate 2 GO (non-branching) / Gate
3 NA}\label{attempt-7-gate-1-go-gate-2-go-non-branching-gate-3-na}}

\begin{itemize}
\tightlist
\item
  \textbf{Held: Science lands (first fully-autonomous success).}
  Correct non-branching read: a fold-free curve rooted at the
  Gast/pSom-high developmental origin, Gate 2 GO on biology, Gate 3 NA.
  The Kürten lessons landed --- curve not tree, root at origin, fold
  checked, no re-embedding.
\item
  \textbf{Fail: Figure style diverges from the binding reference.}
  Panels square, not the house \texttt{(3.5,1.5)} linear-trajectory
  aspect; signature UMAPs plot graph-smoothed scores where the reference
  shows raw (smoothing aids the analysis but misleads as a deliverable);
  EMT-TF marker size doesn't match the other panels.
\item
  \textbf{Fail: \texttt{summary.pdf} clips text; tidiness not fully
  observed.} Verdict and Geometry lines run past the page box, cut off
  mid-sentence.
\item
  \textbf{Watch: One residual island.} A P03-dominated detached lobe
  survives island removal (39-cell speck stripped, not this); not fatal
  to the continuous-manifold call.
\end{itemize}

\begin{center}\rule{0.5\linewidth}{0.5pt}\end{center}

\hypertarget{bischoff-2021-luad-10x}{%
\subsection{Bischoff2021 (LUAD, 10x)}\label{bischoff-2021-luad-10x}}

Lung adenocarcinoma, 10x, \textasciitilde13.4k malignant cells --- new
biology versus the Puram/Kürten/Chen head-and-neck runs.

\hypertarget{attempt-1-gate-1-go-contested-gate-2-go-non-branching-gate-3-na}{%
\subsubsection{Attempt 1 · Gate 1 GO (contested) / Gate 2 GO
(non-branching) / Gate 3
NA}\label{attempt-1-gate-1-go-contested-gate-2-go-non-branching-gate-3-na}}

\begin{itemize}
\tightlist
\item
  \textbf{Fail: Pre-Gate-1 error: two-blob manifold passed as
  continuous.} The manifold splits into two blobs. The agent saw the gap
  but dismissed it as a UMAP artifact because the kNN graph is one
  component --- deciding a visual gate on a connectivity scalar. The
  real cause was never settled: either under-integration (a higher θ,
  rejected for smoothing the gap, may merge the blobs) or a scoping
  failure (identify each lobe's biology and fit the trajectory only on
  the lobe reproducing the reference --- the Kürten evaluation's
  lobe-identity fork).
\item
  \textbf{Fail: \texttt{summary.pdf} rule strikes through the title.}
  The header underline crosses the title descenders (and the Gate-3
  ``exclusivity'' line); otherwise the deliverable matches the reference
  exactly.
\item
  \textbf{Held: Deliverables land.} Correct §1.6 set in the
  correct locations, no scratch, house \texttt{(3.5,1.5)} manifold
  aspect, orientation consistent across UMAPs, raw signature scores
  shown. Binding re-read certified at start and each gate.
\end{itemize}

\hypertarget{attempt-2-gate-1-no-go-gate-2-na-gate-3-na}{%
\subsubsection{Attempt 2 · Gate 1 NO-GO / Gate 2 NA / Gate 3
NA}\label{attempt-2-gate-1-no-go-gate-2-na-gate-3-na}}

\begin{itemize}
\tightlist
\item
  \textbf{Fail: Un-integratable call unsupported by its own sweep.}
  Gate 1 NO-GO (``not continuous'') rests on a plateauing diff-sample
  kNN scalar, yet the sweep's own 16 PC / θ=10 panel shows the gap
  closing into one continuous mass. The mirror of attempt 1 --- scalar
  over embedding --- on the opposite gate.
\item
  \textbf{Fail: Sweep mis-centered.} PCs swept only 12/16/20, and
  \textless20 is not robust (precedents ran 20/40), so the grid sits
  below the robust range. The one panel that did converge is its highest
  PC × highest θ corner, so a properly-centered sweep was never run to
  confirm it.
\item
  \textbf{Fail: \texttt{summary.pdf} layout regressed from the
  reference.} Attempt 1 matched the reference one-pager; this run
  redesigned it into a bespoke layout (colored gate chips, new title, an
  extra lobe-fork block).
\end{itemize}

\hypertarget{attempt-3-gate-1-go-gate-2-go-non-branching-gate-3-na}{%
\subsubsection{Attempt 3 · Gate 1 GO / Gate 2 GO (non-branching) / Gate
3 NA}\label{attempt-3-gate-1-go-gate-2-go-non-branching-gate-3-na}}

\begin{itemize}
\tightlist
\item
  \textbf{Held: Science lands (first Bischoff success).} Correct
  non-branching call, and it settles the Gate-1 question both prior
  attempts got wrong: the two density lobes are poles of one continuous
  manifold, shown on both the embedding and the sweep. Deliverables land
  --- correct §1.6 set, no scratch, reference layout restored.
\item
  \textbf{Fail: \texttt{summary.pdf} step figures low-resolution.}
  The six embedded panels were saved at low DPI, so their in-panel text
  blurs against the sharp vector text of the summary itself.
\end{itemize}

\begin{center}\rule{0.5\linewidth}{0.5pt}\end{center}

\hypertarget{sun-2021-hcc-10x}{%
\subsection{Sun2021 (HCC, 10x)}\label{sun-2021-hcc-10x}}

Hepatocellular carcinoma, 10x, \textasciitilde2.3k malignant cells ---
new biology versus the Puram/Kürten/Chen/Bischoff head-and-neck and lung
runs.

\hypertarget{attempt-1-gate-1-go-gate-2-go-non-branching-gate-3-na}{%
\subsubsection{Attempt 1 · Gate 1 GO / Gate 2 GO (non-branching) / Gate
3 NA}\label{attempt-1-gate-1-go-gate-2-go-non-branching-gate-3-na}}

\begin{itemize}
\tightlist
\item
  \textbf{Held: Science lands (first liver success).} Correct
  non-branching read on new biology: one continuous
  gastrulation→neural-crest axis rooted at the Gast-high pole. Checked
  each lobe's biology before certifying (the Kürten fork), kept the full
  manifold, and self-caught three reviewer findings before re-running.
\item
  \textbf{Fail: EMT-TF panel in the wrong plotting style.} A
  continuous single-gene colormap, not the reference categorical overlay
  of thresholded expressers on a grey manifold; the overlay branch was
  coded but bypassed.
\item
  \textbf{Fail: \texttt{summary.pdf} redesigned, not the reference
  layout.} Bespoke one-pager, not the certified ``GO/NO-GO GATES → STEP
  FIGURES → Parameter choices'' format. Build scratch also left in the
  notebook dir.
\item
  \textbf{BINDING: Verify deliverables against the reference
  visually.} Compare every deliverable to the reference output directly,
  not spec prose alone: each panel must match the reference plotting
  style, and \texttt{summary.pdf} its exact layout. Correct science with
  divergent panels or layout is not a pass.
\end{itemize}

\hypertarget{sec:5}{%
\section{Full run ledger}\label{sec:5}}

Running scorecard --- one row per analyzed dataset, reflecting its
\textbf{last attempt}. Rows are added as datasets are run and finalized
(verdicts entered by the expert). Per-run provenance (cell counts,
params, subtypes, gate mechanics) lives in \texttt{analysis\_log.csv},
joined on \texttt{JoinKey}.

\begin{landscape}
\footnotesize
\setlength{\tabcolsep}{2pt}
\begin{longtable}[]{@{}
  >{\raggedright\arraybackslash}p{(\columnwidth - 24\tabcolsep) * \real{0.081}}
  >{\raggedright\arraybackslash}p{(\columnwidth - 24\tabcolsep) * \real{0.105}}
  >{\raggedright\arraybackslash}p{(\columnwidth - 24\tabcolsep) * \real{0.120}}
  >{\raggedright\arraybackslash}p{(\columnwidth - 24\tabcolsep) * \real{0.240}}
  >{\raggedright\arraybackslash}p{(\columnwidth - 24\tabcolsep) * \real{0.062}}
  >{\raggedright\arraybackslash}p{(\columnwidth - 24\tabcolsep) * \real{0.062}}
  >{\raggedright\arraybackslash}p{(\columnwidth - 24\tabcolsep) * \real{0.058}}
  >{\raggedright\arraybackslash}p{(\columnwidth - 24\tabcolsep) * \real{0.070}}
  >{\raggedright\arraybackslash}p{(\columnwidth - 24\tabcolsep) * \real{0.034}}
  >{\raggedright\arraybackslash}p{(\columnwidth - 24\tabcolsep) * \real{0.034}}
  >{\raggedright\arraybackslash}p{(\columnwidth - 24\tabcolsep) * \real{0.034}}
  >{\raggedright\arraybackslash}p{(\columnwidth - 24\tabcolsep) * \real{0.040}}
  >{\raggedright\arraybackslash}p{(\columnwidth - 24\tabcolsep) * \real{0.060}}@{}}
\toprule\noalign{}
\begin{minipage}[t]{\linewidth}\raggedright
Title
\end{minipage} & \begin{minipage}[t]{\linewidth}\raggedright
Category
\end{minipage} & \begin{minipage}[t]{\linewidth}\raggedright
Disease
\end{minipage} & \begin{minipage}[t]{\linewidth}\raggedright
JoinKey
\end{minipage} & \begin{minipage}[t]{\linewidth}\raggedright
Run start
\end{minipage} & \begin{minipage}[t]{\linewidth}\raggedright
Run end
\end{minipage} & \begin{minipage}[t]{\linewidth}\raggedright
Expert interaction
\end{minipage} & \begin{minipage}[t]{\linewidth}\raggedright
Manifold geometry
\end{minipage} & \begin{minipage}[t]{\linewidth}\raggedright
Gate 1
\end{minipage} & \begin{minipage}[t]{\linewidth}\raggedright
Gate 2
\end{minipage} & \begin{minipage}[t]{\linewidth}\raggedright
Gate 3
\end{minipage} & \begin{minipage}[t]{\linewidth}\raggedright
Hub r
\end{minipage} & \begin{minipage}[t]{\linewidth}\raggedright
Attempts
\end{minipage} \\
\midrule\noalign{}
\endhead
\bottomrule\noalign{}
\endlastfoot
Puram et al.~2017 & Head and Neck & Head and Neck Squamous Cell
Carcinoma & Data\_Puram2017\_Head-and-Neck & --- & --- & Teacher &
Branching & GO & GO & GO & −0.47 & 1 \\
Kürten et al.~2021 & Head and Neck & Head and Neck Squamous Cell
Carcinoma & Data\_Kürten2021\_Head-and-Neck & 2026-07-27 15:08 &
2026-07-27 22:10 & Tutor & Linear & GO & GO & --- & --- & 1 \\
Chen et al.~2020 & Head and Neck & Nasopharyngeal Carcinoma &
Data\_Chen2020\_Head-and-Neck & 2026-07-29 13:19 & 2026-07-29 14:47 &
Reviewer & Linear & GO & GO & --- & --- & 7 \\
Bischoff et al.~2021 & Lung & Lung Adenocarcinoma &
Data\_Bischoff2021\_Lung & 2026-07-29 22:24 & 2026-07-29 23:47 &
Reviewer & Linear & GO & GO & --- & --- & 3 \\
Sun et al.~2021 & Liver/Biliary & Hepatocellular Carcinoma &
Data\_Sun2021\_Liver-Biliary & 2026-07-30 00:05 & 2026-07-30 01:23 &
Reviewer & Linear & GO & GO & --- & --- & 1 \\
Bi et al.~2021 ¹ & Kidney & Renal Cell Carcinoma & Data\_Bi2021\_Kidney
& 2026-07-30 05:04 & 2026-07-30 05:48 & Judge & Linear & GO & GO & --- &
--- & 1 \\
Jansky et al.~2021 ² & Neuroendocrine & Neuroblastoma &
Data\_Jansky2021\_Neuroendocrine & 2026-07-30 13:07 & 2026-07-30 15:08 &
Judge & Linear & GO & GO & --- & --- & 1 \\
Tirosh et al.~2016 ³ & Skin & Melanoma & Data\_Tirosh2016\_Skin &
2026-08-05 22:04 & 2026-08-05 22:49 & Judge & Linear & GO & GO & --- &
--- & 2 \\
Peng et al.~2019 & Pancreas & Pancreatic Ductal Adenocarcinoma &
Data\_Peng2019\_Pancreas & 2026-08-06 21:22 & 2026-08-06 22:52 & Judge &
Linear & GO & GO & --- & --- & 3 \\
Chen et al.~2021 & Prostate & Prostate Cancer & Data\_Chen2021\_Prostate
& 2026-08-05 20:26 & 2026-08-05 21:15 & Judge & Linear & GO & GO & --- &
--- & 2 \\
Lee et al.~2020 ⁴ & Colorectal & Colorectal Carcinoma &
Data\_Lee2020\_Colorectal & 2026-07-30 20:15 & 2026-07-30 21:51 & Judge
& Linear & GO & GO & --- & --- & 1 \\
Wang et al.~2019 & Brain & Glioblastoma & Data\_Wang2019\_Brain &
2026-08-08 15:20 & 2026-08-08 17:28 & Judge & --- & GO & NO-GO & --- &
--- & 8 \\
Wu et al.~2021 ‡ & Breast & Breast Cancer & Data\_Wu2021\_Breast &
2026-07-31 00:32 & 2026-07-31 03:03 & Judge & Linear & GO & GO & --- &
--- & 1 \\
Zhang et al.~2022 & Ovarian & High-Grade Serous Ovarian Carcinoma &
Data\_Zhang2022\_Ovarian & 2026-07-31 03:15 & 2026-07-31 13:56 & Judge &
Linear & GO & GO & --- & --- & 1 \\
Zhou et al.~2020 † & Sarcoma & Osteosarcoma & Data\_Zhou2020\_Sarcoma &
2026-07-31 14:00 & 2026-07-31 16:14 & Judge & Linear & GO & GO & --- &
--- & 1 \\
Kim et al.~2020 & Lung & Lung Adenocarcinoma & Data\_Kim2020\_Lung &
2026-08-06 21:22 & 2026-08-06 21:56 & Judge & Linear & GO & GO & --- &
--- & 3 \\
Steele et al.~2020 & Pancreas & Pancreatic Ductal Adenocarcinoma &
Data\_Steele2020\_Pancreas & 2026-07-31 17:35 & 2026-07-31 18:33 & Judge
& Linear & GO & GO & --- & --- & 1 \\
Song et al.~2022 & Prostate & Prostate Cancer & Data\_Song2022\_Prostate
& 2026-07-31 18:37 & 2026-07-31 19:43 & Judge & Linear & GO & GO & --- &
--- & 1 \\
Ji et al.~2020 & Skin & Skin Squamous Cell Carcinoma &
Data\_Ji2020\_Skin & 2026-07-31 20:16 & 2026-07-31 22:01 & Judge &
Branching & GO & GO & GO & −0.41 & 10 \\
Lin et al.~2020 & Pancreas & Pancreatic Ductal Adenocarcinoma &
Data\_Lin2020\_Pancreas & 2026-08-01 15:13 & 2026-08-01 16:27 & Judge &
Linear & GO & GO & --- & --- & 1 \\
Qian et al.~2020 & Lung & Lung Cancer & Data\_Qian2020\_Lung &
2026-08-01 17:17 & 2026-08-01 18:10 & Judge & Linear & GO & GO & --- &
--- & 1 \\
Obradovic et al.~2021 & Kidney & Renal Cell Carcinoma &
Data\_Obradovic2021\_Kidney & 2026-08-08 15:20 & 2026-08-08 17:51 &
Judge & Linear & GO & GO & --- & --- & 1 \\
Venteicher et al.~2017 & Brain & IDH-Mutant Astrocytoma &
Data\_Venteicher2017\_Brain & 2026-08-02 00:23 & 2026-08-02 01:06 &
Judge & Linear & GO & GO & --- & --- & 1 \\
Olbrecht et al.~2021 & Ovarian & High-Grade Serous Tubo-Ovarian Cancer &
Data\_Olbrecht2021\_Ovarian & 2026-08-08 15:20 & 2026-08-08 19:48 &
Judge & Linear & GO & GO & --- & --- & 2 \\
Geistlinger et al.~2020 & Ovarian & High-Grade Serous Ovarian Carcinoma
& Data\_Geistlinger2020\_Ovarian & 2026-08-05 22:04 & 2026-08-05 22:52 &
Judge & Linear & GO & GO & --- & --- & 2 \\
Laughney et al.~2020 & Lung & Lung Cancer & Data\_Laughney2020\_Lung &
2026-08-02 03:27 & 2026-08-02 15:58 & Judge & Linear & GO & GO & --- &
--- & 1 \\
Zhang et al.~2021 & Kidney & Renal Cell Carcinoma &
Data\_Zhang2021\_Kidney & 2026-08-07 05:53 & 2026-08-07 07:05 & Judge &
Linear & GO & GO & --- & --- & 4 \\
Zilionis et al.~2019 & Lung & Non-Small Cell Lung Cancer &
Data\_Zilionis2019\_Lung & 2026-08-07 23:43 & 2026-08-08 03:09 & Judge &
--- & GO & NO-GO & --- & --- & 6 \\
Choudhury et al.~2022 & Brain & Meningioma & Data\_Choudhury2022\_Brain
& 2026-08-02 18:21 & 2026-08-02 21:46 & Judge & --- & NO-GO & --- & ---
& --- & 1 \\
Dong et al.~2020 & Prostate & Prostate Cancer & Data\_Dong2020\_Prostate
& 2026-08-07 23:43 & 2026-08-08 02:44 & Judge & Linear & GO & GO & --- &
--- & 6 \\
Jerby-Arnon et al.~2018 ⁵ & Skin & Melanoma &
Data\_Jerby-Arnon2018\_Skin & 2026-08-02 23:05 & 2026-08-03 00:08 &
Judge & Linear & GO & GO & --- & --- & 1 \\
Biermann et al.~2022 & Skin & Metastatic Melanoma &
Data\_Biermann2022\_Skin & 2026-08-03 00:52 & 2026-08-03 03:12 & Judge &
Linear & GO & GO & --- & --- & 1 \\
Maynard et al.~2020 & Lung & Lung Cancer & Data\_Maynard2020\_Lung &
2026-08-03 03:33 & 2026-08-03 12:29 & Judge & Linear & GO & GO & --- &
--- & 1 \\
Xing et al.~2021 & Lung & Lung Cancer & Data\_Xing2021\_Lung &
2026-08-03 12:39 & 2026-08-03 13:43 & Judge & Linear & GO & GO & --- &
--- & 1 \\
Chan et al.~2021 ⁶ & Lung & Lung Cancer & Data\_Chan2021\_Lung &
2026-08-07 05:53 & 2026-08-07 07:45 & Judge & Linear & GO & GO & --- &
--- & 4 \\
Galen et al.~2019 & Hematologic & Acute Myeloid Leukemia &
Data\_Galen2019\_Hematologic & 2026-08-03 14:21 & 2026-08-03 15:22 &
Judge & Linear & GO & GO & --- & --- & 1 \\
Gao et al.~2021 & Breast & Breast Cancer & Data\_Gao2021\_Breast &
2026-08-07 05:53 & 2026-08-07 06:44 & Judge & --- & NO-GO & --- & --- &
--- & 4 \\
Hwang et al.~2022 & Pancreas & Pancreatic Ductal Adenocarcinoma &
Data\_Hwang2022\_Pancreas & 2026-08-07 05:53 & 2026-08-07 10:09 & Judge
& --- & GO & NO-GO & --- & --- & 4 \\
Izar et al.~2020 & Ovarian & High-Grade Serous Ovarian Carcinoma &
Data\_Izar2020\_Ovarian & 2026-08-03 14:21 & 2026-08-03 15:03 & Judge &
Linear & GO & GO & --- & --- & 1 \\
Jerby-Arnon et al.~2021 & Sarcoma & Synovial Sarcoma &
Data\_Jerby-Arnon2021\_Sarcoma & 2026-08-05 15:19 & 2026-08-05 16:13 &
Judge & Linear & GO & GO & --- & --- & 2 \\
Kildisiute et al.~2021 & Neuroendocrine & Neuroblastoma &
Data\_Kildisiute2021\_Neuroendocrine & 2026-08-05 15:19 & 2026-08-05
15:52 & Judge & --- & NO-GO & --- & --- & --- & 2 \\
Krishna et al.~2021 & Kidney & Renal Cell Carcinoma &
Data\_Krishna2021\_Kidney & 2026-08-03 14:21 & 2026-08-03 15:03 & Judge
& Linear & GO & GO & --- & --- & 1 \\
Neftel et al.~2019 & Brain & Glioblastoma & Data\_Neftel2019\_Brain &
2026-08-03 14:21 & 2026-08-03 15:25 & Judge & Linear & GO & GO & --- &
--- & 1 \\
Pelka et al.~2021 & Colorectal & Colorectal Cancer &
Data\_Pelka2021\_Colorectal & 2026-08-03 14:21 & 2026-08-03 18:12 &
Judge & Linear & GO & GO & --- & --- & 1 \\
Olalekan et al.~2021 ⁷ & Ovarian & Ovarian Cancer &
Data\_Olalekan2021\_Ovarian & 2026-08-03 19:57 & 2026-08-03 20:21 &
Judge & Linear & GO & GO & --- & --- & 1 \\
Young et al.~2018 & Kidney & Kidney Cancer & Data\_Young2018\_Kidney &
2026-08-05 15:19 & 2026-08-05 15:52 & Judge & --- & NO-GO & --- & --- &
--- & 2 \\
He et al.~2021 & Prostate & Metastatic Prostate Cancer &
Data\_He2021\_Prostate & 2026-08-03 19:57 & 2026-08-04 13:22 & Judge &
Linear & GO & GO & --- & --- & 1 \\
Qian et al.~2020 & Breast & Breast Cancer & Data\_Qian2020\_Breast &
2026-08-06 21:22 & 2026-08-06 22:36 & Judge & --- & NO-GO & --- & --- &
--- & 3 \\
Paulson et al.~2020 & Skin & Merkel Cell Carcinoma &
Data\_Paulson2020\_Skin & 2026-08-03 19:57 & 2026-08-03 20:31 & Judge &
Linear & GO & GO & --- & --- & 1 \\
Yuan et al.~2018 & Brain & High-Grade Glioma & Data\_Yuan2018\_Brain &
2026-08-03 19:57 & 2026-08-03 20:29 & Judge & Linear & GO & GO & --- &
--- & 1 \\
Gojo et al.~2020 & Brain & Pediatric Ependymoma & Data\_Gojo2020\_Brain
& 2026-08-03 19:57 & 2026-08-04 13:22 & Judge & Linear & GO & GO & --- &
--- & 1 \\
Raghavan et al.~2021 & Pancreas & Pancreatic Ductal Adenocarcinoma &
Data\_Raghavan2021\_Pancreas & 2026-08-07 18:53 & 2026-08-07 19:41 &
Judge & Linear & GO & GO & --- & --- & 5 \\
Cohen et al.~2021 & Hematologic & Multiple Myeloma &
Data\_Cohen2021\_Hematologic & 2026-08-06 21:22 & 2026-08-06 23:17 &
Judge & Linear & GO & GO & --- & --- & 3 \\
Dong et al.~2020 & Neuroendocrine & Neuroblastoma &
Data\_Dong2020\_Neuroendocrine & 2026-08-03 19:57 & 2026-08-03 21:10 &
Judge & Linear & GO & GO & --- & --- & 1 \\
Couturier et al.~2020⁸ & Brain & Glioblastoma &
Data\_Couturier2020\_Brain & 2026-08-04 14:02 & 2026-08-04 15:51 & Judge
& Linear & GO & GO & --- & --- & 1 \\
Filbin et al.~2018 & Brain & Pediatric H3-K27M Mutant Glioma &
Data\_Filbin2018\_Brain & 2026-08-04 17:04 & 2026-08-04 17:12 & Judge &
Linear & GO & GO & --- & --- & 1 \\
Gaiti et al.~2019 & Hematologic & Chronic Lymphocytic Leukemia &
Data\_Gaiti2019\_Hematologic & 2026-08-04 14:02 & 2026-08-04 15:29 &
Judge & Linear & GO & GO & --- & --- & 1 \\
Liu et al.~2021 & Hematologic & Multiple Myeloma &
Data\_Liu2021\_Hematologic & 2026-08-05 15:19 & 2026-08-05 16:20 & Judge
& --- & NO-GO & --- & --- & --- & 2 \\
Nath et al.~2021 & Ovarian & High-Grade Serous Ovarian Carcinoma &
Data\_Nath2021\_Ovarian & 2026-08-04 14:02 & 2026-08-04 15:36 & Judge &
Linear & GO & GO & --- & --- & 1 \\
Pal et al.~2021 & Breast & Breast Cancer & Data\_Pal2021\_Breast &
2026-08-06 21:22 & 2026-08-07 04:24 & Judge & Linear & GO & GO & --- &
--- & 3 \\
Zhang et al.~2019 & Ovarian & High-Grade Serous Ovarian Carcinoma &
Data\_Zhang2019\_Ovarian & 2026-08-04 14:02 & 2026-08-04 15:07 & Judge &
Linear & GO & GO & --- & --- & 1 \\
Li et al.~2017 & Colorectal & Colorectal Carcinoma &
Data\_Li2017\_Colorectal & 2026-08-04 15:54 & 2026-08-04 16:20 & Judge &
Linear & GO & GO & --- & --- & 1 \\
Tirosh et al.~2016 & Brain & IDH-Mutant Oligodendroglioma &
Data\_Tirosh2016\_Brain & 2026-08-04 15:54 & 2026-08-04 16:20 & Judge &
Linear & GO & GO & --- & --- & 1 \\
Bassez et al.~2021 & Breast & Breast Cancer & Data\_Bassez2021\_Breast &
2026-08-07 23:43 & 2026-08-08 03:19 & Judge & Linear & GO & GO & --- &
--- & 6 \\
Caron et al.~2020 & Hematologic & Acute Lymphoblastic Leukemia &
Data\_Caron2020\_Hematologic & 2026-08-07 06:57 & 2026-08-07 10:40 &
Judge & Linear & GO & GO & --- & --- & 4 \\
Griffiths et al.~2021 & Breast & Breast Cancer &
Data\_Griffiths2021\_Breast & 2026-08-07 23:43 & 2026-08-08 04:38 &
Judge & --- & GO & NO-GO & --- & --- & 6 \\
Hovestadt et al.~2019 & Brain & Medulloblastoma &
Data\_Hovestadt2019\_Brain & 2026-08-04 16:41 & 2026-08-04 17:16 & Judge
& Linear & GO & GO & --- & --- & 1 \\
Ledergor et al.~2018 & Hematologic & Multiple Myeloma &
Data\_Ledergor2018\_Hematologic & 2026-08-05 16:54 & 2026-08-05 17:49 &
Judge & Linear & GO & GO & --- & --- & 2 \\
Nam et al.~2019 & Hematologic & Myeloproliferative Neoplasms &
Data\_Nam2019\_Hematologic & 2026-08-04 16:41 & 2026-08-04 17:36 & Judge
& Linear & GO & GO & --- & --- & 1 \\
Qian et al.~2020 & Ovarian & Ovarian Cancer & Data\_Qian2020\_Ovarian &
2026-08-05 16:54 & 2026-08-05 17:28 & Judge & Linear & GO & GO & --- &
--- & 2 \\
Rendeiro et al.~2020 & Hematologic & Chronic Lymphocytic Leukemia &
Data\_Rendeiro2020\_Hematologic & 2026-08-07 23:43 & 2026-08-08 02:43 &
Judge & --- & NO-GO & --- & --- & --- & 6 \\
Roider et al.~2020 & Hematologic & B-Cell Lymphoma &
Data\_Roider2020\_Hematologic & 2026-08-04 16:41 & 2026-08-04 17:21 &
Judge & Linear & GO & GO & --- & --- & 1 \\
Wu et al.~2020 & Hematologic & Acute Myeloid Leukemia &
Data\_Wu2020\_Hematologic & 2026-08-07 07:47 & 2026-08-07 10:40 & Judge
& --- & GO & NO-GO & --- & --- & 4 \\
Chung et al.~2017 & Breast & Breast Cancer & Data\_Chung2017\_Breast &
2026-08-04 22:59 & 2026-08-05 01:42 & Judge & Linear & GO & GO & --- &
--- & 1 \\
Darmanis et al.~2017 & Brain & Glioblastoma & Data\_Darmanis2017\_Brain
& 2026-08-05 15:47 & 2026-08-05 15:56 & Judge & Linear & GO & GO & --- &
--- & 2 \\
Giustacchini et al.~2017 & Hematologic & Chronic Myeloid Leukemia &
Data\_Giustacchini2017\_Hematologic & 2026-08-05 20:26 & 2026-08-05
20:49 & Judge & Linear & GO & GO & --- & --- & 2 \\
Gulati et al.~2020 & Breast & Breast Cancer & Data\_Gulati2020\_Breast &
2026-08-05 22:04 & 2026-08-05 22:25 & Judge & --- & GO & NO-GO & --- &
--- & 2 \\
Karaayvas et al.~2018 & Breast & Breast Cancer &
Data\_Karaayvas2018\_Breast & 2026-08-05 15:47 & 2026-08-05 16:00 &
Judge & --- & GO & NO-GO & --- & --- & 2 \\
Kim et al.~2018 & Breast & Triple-Negative Breast Cancer &
Data\_Kim2018\_Breast & 2026-08-04 22:23 & 2026-08-04 23:39 & Judge &
Linear & GO & GO & --- & --- & 1 \\
Li et al.~2019 & Skin & Melanoma & Data\_Li2019\_Skin & 2026-08-07 09:21
& 2026-08-07 10:40 & Judge & --- & NO-GO & --- & --- & --- & 4 \\
Moncada et al.~2020 & Pancreas & Pancreatic Ductal Adenocarcinoma &
Data\_Moncada2020\_Pancreas & 2026-08-04 22:55 & 2026-08-05 01:42 &
Judge & Linear & GO & GO & --- & --- & 1 \\
Shih et al.~2018 & Ovarian & Ovarian Cancer & Data\_Shih2018\_Ovarian &
2026-08-04 22:23 & 2026-08-04 23:15 & Judge & Linear & GO & GO & --- &
--- & 1 \\
Song et al.~2019 & Lung & Non-Small Cell Lung Cancer &
Data\_Song2019\_Lung & 2026-08-04 22:23 & 2026-08-04 23:04 & Judge &
Linear & GO & GO & --- & --- & 1 \\
Steen et al.~2021 & Hematologic & Diffuse Large B-Cell Lymphoma &
Data\_Steen2021\_Hematologic & 2026-08-04 22:23 & 2026-08-05 01:11 &
Judge & Linear & GO & GO & --- & --- & 1 \\
Yost et al.~2019 & Skin & Skin Carcinoma & Data\_Yost2019\_Skin &
2026-08-04 22:23 & 2026-08-04 22:54 & Judge & Linear & GO & GO & --- &
--- & 1 \\
Zhang et al.~2019 & Hematologic & Follicular Lymphoma &
Data\_Zhang2019\_Hematologic & 2026-08-07 23:43 & 2026-08-08 02:19 &
Judge & --- & GO & NO-GO & --- & --- & 6 \\
\end{longtable}
\end{landscape}

\textbf{Legend.}

\begin{itemize}
\item
  \textbf{Expert interaction}: Teacher (calibrated the protocol) ·
  Tutor (steered step-by-step) · Reviewer (autonomous run, audited after
  the fact) · Judge (outcome adjudicated, no steering).
\item
  \textbf{Manifold geometry}: Branching (Gast/pSom hub → pSom arm and Gast arm, cross-repression; Crest a terminal readout at the arm tips)
  · Linear (collapsed bifurcation, single continuous axis).
\item
  \textbf{Gate 1 · Gate 2 · Gate 3}: GO / NO-GO. Gate 3 is
  branching-only.
\item
  \textbf{Hub r}: Gast↔pSom cross-repression Pearson; branching only.
\item
  \textbf{Attempts}: Number of runs to reach this finalized row.
\item
  \textbf{‡ Wu et al.~2021}: The agent identified an ambiguity in the
  non-branching definition, which did not say whether Gast/pSom
  cross-repression is disqualifying. The expert then updated the
  definition (§1.1--2) to include the cross-repression case ---
  organized cross-repression with a terminal Crest gradient counts as
  non-branching. Agent-identified, expert-ratified.
\item
  \textbf{† Zhou et al.~2020}: Both this run and Zhang 2022 struggled
  to keep orientation consistent across UMAPs, so the expert made it a
  binding criterion (Section 4), which this run then met.
\end{itemize}

\textbf{Scope corrections.} \textbf{Ireland et al.~2020 (Lung, SCLC)}
--- removed (row, both \texttt{analysis\_log.csv} rows and all deliverables), dropping the atlas from 87 to 86 rows. It is a
cross-species composite (mouse RPM GEM time-course and human sample), so
it is not validly in scope.

\textbf{Prompt Evolution.} The autonomous \textbf{Judge} runs shared one
initiating prompt that set only the regime --- scientific instruction
stayed in the spec files (except the branching re-run, item 12, which
carried scientific steering in the prompt and via debrief). Its
substantive edits are numbered below; each carries a superscript on the
table row where it first appears (re-run waves overwrite existing rows,
so carry none). Operational guardrails (do-not-interrupt, no new
packages/environments) are omitted.

\begin{enumerate}
\def\labelenumi{\arabic{enumi}.}
\tightlist
\item
  \textbf{Baseline regime}: Named dataset; no-feedback mandate;
  post-hoc independent-auditor contract. The seed paragraph, carried
  verbatim thereafter.
\item
  \textbf{Agent-chosen dataset}: Dataset choice delegated to the
  agent; deliverables first named as the recurring weak spot.
\item
  \textbf{Learn from the log}: Reuse deliverable templates from prior
  \emph{successful} runs, especially \texttt{summary.pdf}.
\item
  \textbf{Science and deliverables}: The flagged weak spot widened from
  deliverables to include science; this became the canonical template
  for the single-dataset runs that followed.
\item
  \textbf{Content steer} (one run only): The sole nudge on
  \emph{what} to analyze: pick a low-pSom dataset. The next run reverted
  to item 4, unmarked.
\item
  \textbf{Batch delegation / scale-up}: The item-4 template plus one
  instruction: select 10 datasets and delegate the identical autonomous
  instructions to 10 subagents, reviewed together. First multi-dataset
  run.
\item
  \textbf{Delegate self-check}: Item 6 plus one clause: the
  delegating agent may also check its subagents' science and deliverables
  accuracy. Second batch.
\item
  \textbf{Full scale-up}: Item 7 unchanged in regime, scope widened
  from ``10 new'' to ``all remaining in-scope datasets'': the final
  batch, 33 subagents delegated the identical instructions and reviewed
  together.
\item
  \textbf{Science-fail re-run (\texttt{-1})}: Item-8 regime plus one
  directive: re-analyze every first-run science failure, hinting most
  were bad Gate 1/2 GO calls or missed branching.
\item
  \textbf{Science-fail re-run (\texttt{-3})}: Item-9 regime
  re-applied to the datasets still failing science. Dispatch
  instructions added strict attention to Gate 1 integration and Gate 2
  module geography, cycling-lobe scoping, pseudotime-curve rooting and
  fold-freeness, and deliverable/biology-visualization fidelity.
\item
  \textbf{Science-fail re-run (\texttt{-5})}: Item-10 regime,
  dispatch retargeted from pseudotime rooting/fold-freeness to which
  region the trajectory is scoped on. Several Gate 2 GOs flipped to
  NO-GO. This exposed a gap in the protocol order: \texttt{02\_manifold}
  fits the pseudotime before Gate 2 tests the modules, so the fitted
  trajectory does not depend on the embryonic modules being present.
\item
  \textbf{Targeted content steer}: For the last run of the
  bifurcating dataset (Ji), the item-11 regime was narrowed with
  scientific steering toward a branching outcome (in-prompt and via
  debrief; see header).
\end{enumerate}
\end{document}
